
\section*{Ex.1 $|x^2+1|<x+1$}
Per poder llevar el valor absolut passam:
\begin{equation*}
|x^2+1|<x+1 \ \Longleftrightarrow \ x^2+1<x+1 \ \wedge \ x^2+1>-x-1
\end{equation*}
\\
Començam per la primera part:
\\
\begin{equation*}
x^2+1<x+1 \ \Leftrightarrow \ x^2-x-2<0 
\end{equation*}
\\
Passam la ineqüació a una equació de segon grau: $x^2-x-2=0$ que té per solucions $x=2 \ \wedge \ x=-1$. \\
\\
Cream una taula pels valors:
\begin{table}[h!]
\begin{center}
\begin{tabular}{c|c|c}
$\oplus$ & $\ominus$ & $\oplus$\\
\hline
$(-\infty,-1)$ & $(-1,2)$ & $(2,+\infty)$
\end{tabular}
\end{center}
\end{table}
\\
Per tant la primera solució és $(-1,2)$.
\\
\\
Ens centram ara en la segona part:
\begin{equation*}
x^2+1>-x-1 \ \Leftrightarrow \ x^2+x>0 \ \Leftrightarrow \ x(x+1)>0
\end{equation*}
Ara resolem la casuítica:
\begin{equation*}
\left.
\begin{array}{cc}
\left\lbrace
\begin{array}{cc}
x>0\\
x+1>0
\end{array}\right.
\ \Longrightarrow \
\left\lbrace
\begin{array}{cc}
x>0\\
x>-1
\end{array}\right.
\ \Longrightarrow \
x\in (0,+\infty)
\\
\\
\left\lbrace
\begin{array}{cc}
x<0\\
x+1<0
\end{array}\right.
\ \Longrightarrow \
\left\lbrace
\begin{array}{cc}
x<0\\
x<-1
\end{array}\right.
\ \Longrightarrow \
x\in(-\infty,-1)
\end{array}\right\rbrace
\ \Longrightarrow \
x\in(-\infty,-1)\cup(0,+\infty)
\end{equation*}
Per tant la segona solució és $(-\infty,-1)\cup(0,+\infty)$.
\\
\\
\\
La solució de la inequació és la intersseció entre les dues solucions, és a dir:
\begin{equation*}
(-1,2)\cap((-\infty,-1)\cup(0,+\infty))=(0,2) \ \Longrightarrow \ x\in(0,2)
\end{equation*}
\newpage
\section*{Ex.2 $x^3-x^2+2x-2<0$}
Primer, la passam a una equació de tercer grau com $x^3-x^2+2x-2=0$ que té per solució real $x=1$ i dues solucions imaginàries que són prescindibles com són $x=\sqrt{2}i \ \wedge \ x=-\sqrt{2}i$.
\\
Cream una taula de valors:
\begin{table}[h!]
\begin{center}
\begin{tabular}{c|c}
$\ominus$ & $\oplus$\\
\hline
$(-\infty,1)$ & $(1,+\infty)$
\end{tabular}
\end{center}
\end{table}
\\
Per tant la solució de la inequació és:
$(-\infty, 1) \ \Longrightarrow \ x\in(-\infty, 1)$
\\
\section*{Ex.3 $|x+3|+|x-1|\geq5$}
Primer desenvolupam els valors absoluts:
\begin{equation*}
|x+3|=\left\lbrace
\begin{array}{cc}
x+3 \ \ \ \text{si} \ \ x\geq-3
\\ -x-3 \ \ \ \text{si} \ \ x<-3
\end{array}\right.
\ \wedge \
|x-1|=\left\lbrace
\begin{array}{cc}
x-1 \ \ \ \text{si} \ \ x\geq1
\\ -x+1 \ \ \ \text{si} \ \ x<1
\end{array}\right.
\end{equation*}
Ara resolem la casuística per trams que venen mercats pels punts de canvi en els valors absoluts:
\begin{itemize}
\item $x<-3$:
\begin{equation*}
|x+3|+|x-1|\geq5 \Leftrightarrow -x-3-x+1\geq5 \Leftrightarrow -2x-2\geq5 \Leftrightarrow -2x\geq7 \Leftrightarrow x\leq-\frac{7}{2} \ \Longrightarrow \ x\in\Bigg(-\infty,-\frac{7}{2}\Bigg]
\end{equation*}
\item $x\in[-3,1)$:
\begin{equation*}
|x+3|+|x-1|\geq5 \Leftrightarrow x+3-x+1\geq5 \Leftrightarrow 4\geq5
\end{equation*}
Però $4\geq5$ no és cert perquè $4\notin[-3,1) \ \Longrightarrow$ la solució és $\emptyset$.
\item $x\geq1$:
\begin{equation*}
|x+3|+|x-1|\geq5 \Leftrightarrow x+3+x-1\geq5 \Leftrightarrow 2x+2\geq5 \Leftrightarrow 2x\geq3 \Leftrightarrow x\leq\frac{3}{2} \ \Longrightarrow \ x\in\Bigg[\frac{3}{2},+\infty\Bigg)
\end{equation*}
\end{itemize}
En definitva, la solució a la inequació és:
\begin{equation*}
\Bigg(-\infty,-\frac{7}{2}\Bigg]\cup\emptyset\cup\Bigg[\frac{3}{2},+\infty\Bigg) = \Bigg(-\infty,-\frac{7}{2}\Bigg]\cup\Bigg[\frac{3}{2},+\infty\Bigg)=\mathbb{R}\setminus\Big(-\frac{7}{2},\frac{3}{2}\Big)
\end{equation*}